% =============================================================================
%  Fundamentos compartilhados --- Módulo 05: interfaces declarativas
%  Incluído pelas trilhas Kotlin e Flutter com % =============================================================================
%  Fundamentos compartilhados --- Módulo 05: interfaces declarativas
%  Incluído pelas trilhas Kotlin e Flutter com % =============================================================================
%  Fundamentos compartilhados --- Módulo 05: interfaces declarativas
%  Incluído pelas trilhas Kotlin e Flutter com % =============================================================================
%  Fundamentos compartilhados --- Módulo 05: interfaces declarativas
%  Incluído pelas trilhas Kotlin e Flutter com \input{../fundamentos/f05-ui-declarativa}
% =============================================================================
\section{Fundamentos: interfaces declarativas}

\begin{frame}{Interfaces imperativas e declarativas}
  \begin{columns}[T]
    \begin{column}{0.49\textwidth}
      \begin{definicao}[Interface imperativa]
        O programa cria os componentes uma vez e depois os \textbf{altera} passo a passo (\cod{setText}, \cod{setVisibility}). Basta esquecer um deles para a tela mostrar dados errados.
      \end{definicao}
    \end{column}
    \begin{column}{0.49\textwidth}
      \begin{definicao}[Interface declarativa]
        O programa \textbf{descreve} como a tela deve ser para cada estado. Quando o estado muda, o framework gera a nova descrição e aplica apenas as diferenças.
      \end{definicao}
    \end{column}
  \end{columns}
  \vspace{2pt}
  \centering
  \begin{tikzpicture}[node distance=2mm and 6mm,every node/.style={font=\scriptsize}]
    \node[rotulo,anchor=east,text width=1.6cm,align=right] (li) at (0,0) {\textbf{imperativo}};
    \node[caixa,right=3mm of li,minimum width=1.6cm] (e1) {estado muda};
    \node[caixaambar,right=of e1,minimum width=3.2cm] (c1) {seu código acha e altera\\ cada componente afetado};
    \node[caixa,right=of c1,minimum width=2.2cm] (t1) {tela (mutável)};
    \draw[seta] (e1) -- (c1); \draw[seta] (c1) -- (t1);
    \node[rotulo,anchor=east,text width=1.6cm,align=right] (ld) at (0,-1.05) {\textbf{declarativo}};
    \node[caixa,right=3mm of ld,minimum width=1.6cm] (e2) {estado muda};
    \node[caixapetroleo,right=of e2,minimum width=2.1cm] (f2) {UI = f(estado)};
    \node[caixa,right=of f2,minimum width=2.1cm] (d2) {nova descrição};
    \node[caixapetroleo,right=of d2,minimum width=2.6cm] (a2) {framework compara\\ e aplica diferenças};
    \draw[seta] (e2) -- (f2); \draw[seta] (f2) -- (d2); \draw[seta] (d2) -- (a2);
  \end{tikzpicture}
  \raggedright
  \begin{historia}
    Views do Android e UIKit do iOS (2008) são imperativos. O React (2013) popularizou a UI declarativa, depois adotada pelo Flutter (2018), SwiftUI (2019) e Jetpack Compose (2021).
  \end{historia}
\end{frame}

\begin{frame}{Composição: a tela como uma árvore}
  \begin{columns}[c]
    \begin{column}{0.52\textwidth}
      \centering
      \begin{tikzpicture}[level distance=10mm,
          every node/.style={caixa,font=\tiny,minimum height=5mm,inner sep=2pt},
          edge from parent/.style={draw=azulpetroleo,line width=0.6pt},
          level 1/.style={sibling distance=30mm},
          level 2/.style={sibling distance=21mm},
          level 3/.style={sibling distance=10mm}]
        \node[caixaescura,font=\tiny] (raiz) {TelaEquipamentos}
          child {node {BarraSuperior}}
          child {node {ListaEquipamentos}
            child {node[caixaambar,font=\tiny] {CartaoBancada}
              child {node {Ícone}}
              child {node {Textos}}
              child {node {Barra}}}
            child {node[caixaambar,font=\tiny] {CartaoBancada}}};
        \draw[seta] (-2.6,0.2) -- (-2.6,-3.2);
        \draw[setaambar] (-2.2,-3.2) -- (-2.2,0.2);
        \node[rotulo,anchor=north west,text width=6cm,align=left] at (-2.75,-3.45) {\textcolor{azulpetroleo}{$\downarrow$} dados descem por parâmetros\quad \textcolor{ambar}{$\uparrow$} eventos sobem por \emph{callbacks}};
      \end{tikzpicture}
    \end{column}
    \begin{column}{0.46\textwidth}
      \begin{definicao}[Composição]
        Construir interfaces complexas \textbf{combinando} componentes pequenos e independentes, em vez de herdar de componentes grandes. Cada componente recebe dados e devolve uma descrição de UI.
      \end{definicao}
      \begin{analogia}
        Como peças de montar: o cartão é feito de ícone, textos e indicador; a lista, de vários cartões. Trocar uma peça não obriga a refazer as outras.
      \end{analogia}
      \begin{dica}[Um bom componente]
        Faz uma coisa só, recebe tudo por parâmetro, não sabe de onde vêm os dados e pode ser visto isoladamente em uma \emph{preview}.
      \end{dica}
    \end{column}
  \end{columns}
\end{frame}

\begin{frame}{Layout por restrições}
  \begin{columns}[c]
    \begin{column}{0.47\textwidth}
      \centering
      \begin{tikzpicture}[node distance=11mm]
        \node[caixaescura,minimum width=3.4cm] (t) {Tela \quad 360 $\times$ 800};
        \node[caixa,minimum width=3.4cm,below=of t] (c) {Coluna};
        \node[caixaambar,minimum width=3.4cm,below=of c] (k) {Cartão};
        \draw[seta] ([xshift=-10mm]t.south) -- node[rotulo,left,text width=2.1cm,align=right]{largura 0--360\\ altura 0--800} ([xshift=-10mm]c.north);
        \draw[seta] ([xshift=-10mm]c.south) -- node[rotulo,left,text width=2.1cm,align=right]{largura 0--360\\ altura livre ($\infty$)} ([xshift=-10mm]k.north);
        \draw[setaambar] ([xshift=10mm]k.north) -- node[rotulo,right,text width=1.6cm,align=left]{``tenho\\ 360 $\times$ 96''} ([xshift=10mm]c.south);
        \node[rotulo,below=1.5mm of k,text width=4cm] {a coluna põe o cartão em $y = 0$\\ e o próximo em $y = 96$};
      \end{tikzpicture}
    \end{column}
    \begin{column}{0.51\textwidth}
      \begin{definicao}[Restrições (\emph{constraints})]
        Intervalo de tamanhos que o pai aceita para o filho: largura e altura mínimas e máximas. \textbf{Justa} quando mínimo = máximo; \textbf{frouxa} quando o mínimo é zero.
      \end{definicao}
      \begin{conceito}[O algoritmo em três frases]
        \textbf{Restrições descem}: o pai diz os limites. \textbf{Tamanhos sobem}: o filho escolhe um tamanho dentro deles. \textbf{O pai posiciona}: o filho não sabe onde está.
      \end{conceito}
      \begin{pordentro}[Uma passada só]
        Cada nó é medido uma única vez por quadro, o que mantém o layout rápido mesmo com milhares de nós. No Compose, cada \cod{Modifier} ajusta as restrições antes de repassá-las; no Flutter, \cod{Padding} e \cod{SizedBox} fazem esse papel.
      \end{pordentro}
    \end{column}
  \end{columns}
\end{frame}

\begin{frame}{Do estado aos pixels: o caminho de um quadro}
  \centering
  \begin{tikzpicture}[node distance=5mm,every node/.style={font=\scriptsize}]
    \node[caixaambar,minimum width=2cm,minimum height=11mm] (s) {estado\\ muda};
    \node[caixapetroleo,right=of s,minimum width=2.3cm,minimum height=11mm] (c) {\textbf{composição}\\ \tiny o que mostrar};
    \node[caixapetroleo,right=of c,minimum width=2.3cm,minimum height=11mm] (l) {\textbf{layout}\\ \tiny tamanho e posição};
    \node[caixapetroleo,right=of l,minimum width=2.3cm,minimum height=11mm] (d) {\textbf{desenho}\\ \tiny cores, textos, formas};
    \node[caixaescura,right=of d,minimum width=2.3cm,minimum height=11mm] (r) {\textbf{rasterização}\\ \tiny GPU gera os pixels};
    \draw[seta] (s) -- (c); \draw[seta] (c) -- (l); \draw[seta] (l) -- (d); \draw[seta] (d) -- (r);
    \node[rotulo,below=1mm of c] {Compose: \emph{Composition}\\ Flutter: \emph{build}};
    \node[rotulo,below=1mm of l] {\emph{Layout}\\ \emph{layout}};
    \node[rotulo,below=1mm of d] {\emph{Drawing}\\ \emph{paint} e \emph{composite}};
    \node[rotulo,below=1mm of r] {motor gráfico\\ (Skia, Impeller)};
  \end{tikzpicture}
  \raggedright
  \begin{columns}[T]
    \begin{column}{0.49\textwidth}
      \begin{definicao}[Quadro e orçamento de tempo]
        A tela é redesenhada no ritmo do sinal vertical do display (VSync): a 60\,Hz cada quadro tem \textbf{16,7\,ms}; a 120\,Hz, só \textbf{8,3\,ms}. Se o trabalho não couber, o quadro se perde e a animação engasga (\emph{jank}).
      \end{definicao}
    \end{column}
    \begin{column}{0.49\textwidth}
      \begin{pordentro}[Pular fases economiza tempo]
        Uma cor que anima só precisa de desenho; uma posição que muda, de layout e desenho. Ler o estado na fase mais tardia possível evita repetir as anteriores.
      \end{pordentro}
      \begin{dica}
        O \emph{Performance} do Flutter DevTools e o \emph{Layout Inspector} do Android Studio mostram onde o tempo do quadro foi gasto.
      \end{dica}
    \end{column}
  \end{columns}
\end{frame}

\begin{frame}{Por dentro: descrição, instância e renderização}
  \begin{columns}[T]
    \begin{column}{0.56\textwidth}
      \scriptsize
      \begin{tabularx}{\linewidth}{@{}L l l@{}}
        \toprule
        \cabecalho{Papel} & \cabecalho{Flutter} & \cabecalho{Compose} \\
        \midrule
        Descrição leve, recriada a cada mudança & \cod{Widget} & chamada \cod{@Composable} \\
        Instância duradoura, guarda o estado & \cod{Element} + \cod{State} & composição (\cod{remember}) \\
        Mede, posiciona e desenha & \cod{RenderObject} & \cod{LayoutNode} \\
        Informa leitores de tela & árvore de semântica & árvore de semântica \\
        \bottomrule
      \end{tabularx}
      \normalsize
      \begin{analogia}
        O widget é a planta do arquiteto: barata de redesenhar. O elemento é o apartamento construído, onde mora o estado. O objeto de renderização é a estrutura medida e pintada. Mudar a planta não obriga a demolir o prédio.
      \end{analogia}
    \end{column}
    \begin{column}{0.42\textwidth}
      \begin{conceito}[Por que separar]
        Descrições são criadas aos milhares por segundo, então precisam ser baratas. Medir e desenhar é caro, então os objetos de renderização são \textbf{reaproveitados} enquanto possível.
      \end{conceito}
      \begin{pordentro}[Semântica e acessibilidade]
        Em paralelo, o framework monta uma árvore de semântica (rótulos, papéis, ações) que o TalkBack e o VoiceOver leem. Ícone sem rótulo fica mudo para quem não enxerga (Módulo 04).
      \end{pordentro}
    \end{column}
  \end{columns}
\end{frame}

\begin{frame}{Atualizar sem refazer tudo}
  \begin{columns}[c]
    \begin{column}{0.42\textwidth}
      \centering
      \begin{tikzpicture}[level distance=10mm,
          every node/.style={caixa,font=\tiny,minimum height=5mm,minimum width=9mm,inner sep=2pt},
          edge from parent/.style={draw=azulpetroleo,line width=0.6pt},
          level 1/.style={sibling distance=22mm},
          level 2/.style={sibling distance=11mm}]
        \node {Tela}
          child {node {Barra} child {node {Título}}}
          child {node[caixaambarforte] (m) {Medidor}
            child {node[caixaambar] {Valor}}
            child {node[caixaambar] {Barra}}};
        \node[rotulo,below=11mm of m,text width=3.4cm] {o estado \cod{temperatura} mudou: só o \textbf{Medidor} e seus filhos são reconstruídos};
      \end{tikzpicture}
    \end{column}
    \begin{column}{0.56\textwidth}
      \begin{conceito}[Flutter: \emph{rebuild} e reconciliação]
        \cod{setState} marca o elemento como sujo; no próximo quadro, só ele executa \cod{build}. Filhos de mesmo tipo e chave são reaproveitados e só atualizados.
      \end{conceito}
      \begin{conceito}[Compose: recomposição inteligente]
        O Compose sabe quais funções leram cada estado e só executa essas de novo; as que recebem parâmetros iguais aos anteriores são puladas (\emph{skipping}).
      \end{conceito}
      \begin{atencao}[Consequência importante]
        Funções de interface podem rodar muitas vezes, em qualquer ordem: devem ser \textbf{rápidas e sem efeitos colaterais}. Rede e banco ficam em APIs de efeito (Módulo 07).
      \end{atencao}
    \end{column}
  \end{columns}
\end{frame}

\begin{frame}[fragile]{Identidade e chaves}
  \begin{columns}[T]
    \begin{column}{0.47\textwidth}
      \centering
      \begin{tikzpicture}[baseline=(current bounding box.north),every node/.style={font=\tiny}]
        \tikzset{item/.style={caixa,minimum width=1.5cm,minimum height=5mm,inner sep=1pt}}
        \node[rotulo] at (0.75,0.55) {antes};
        \node[item,caixaambar] (a) at (0.75,0) {A \quad expandido};
        \node[item] at (0.75,-0.6) {B};
        \node[item] at (0.75,-1.2) {C};
        \node[rotulo,text width=2cm] at (2.95,0.55) {remove A, \textbf{sem} chave};
        \node[item,caixaambar,draw=vermelhoerro] at (2.95,0) {B \quad expandido};
        \node[item] at (2.95,-0.6) {C};
        \node[rotulo,text width=2cm] at (5.15,0.55) {remove A, \textbf{com} chave};
        \node[item] at (5.15,0) {B};
        \node[item] at (5.15,-0.6) {C};
        \node[rotulo,text width=2cm,text=vermelhoerro] at (2.95,-1.25) {o estado ficou na posição 1};
        \node[rotulo,text width=2cm,text=verdeok] at (5.15,-1.25) {o estado saiu junto com A};
      \end{tikzpicture}
\begin{codigo}[title=Compose]{Kotlin}
items(itens, key = { it.id }) { eq ->
  ItemEquipamento(eq)
}
\end{codigo}
\begin{codigo}[title=Flutter]{Dart}
ItemEquipamento(
    key: ValueKey(eq.id), eq: eq, onTap: abrir)
\end{codigo}
    \end{column}
    \begin{column}{0.51\textwidth}
      \begin{definicao}[Chave (\emph{key})]
        Identificador estável que diz ao framework qual item é qual entre duas versões da árvore. Sem chave, a correspondência é feita pela \textbf{posição}.
      \end{definicao}
      \begin{conceito}[Quando importa]
        Em listas cujos itens têm estado próprio (expandido, campo digitado, animação) e podem ser inseridos, removidos ou reordenados.
      \end{conceito}
      \begin{dica}
        Use um identificador do domínio, como o \cod{id} do equipamento. Nunca o índice da lista nem um valor aleatório criado a cada reconstrução.
      \end{dica}
    \end{column}
  \end{columns}
\end{frame}

\begin{frame}{Temas, \emph{tokens} e contexto implícito}
  \begin{columns}[c]
    \begin{column}{0.42\textwidth}
      \centering
      \begin{tikzpicture}[node distance=7mm and 3mm,every node/.style={font=\tiny}]
        \node[caixaambarforte,minimum width=2.4cm,font=\tiny\bfseries] (t) {Tema};
        \node[caixa,below=of t,minimum width=1.6cm,font=\tiny] (s) {Tela};
        \node[caixa,below left=of s,minimum width=1.3cm,font=\tiny] (l) {Lista};
        \node[caixa,below right=of s,minimum width=1.3cm,font=\tiny] (b) {Botão};
        \node[caixa,below=of l,minimum width=1.3cm,font=\tiny] (c) {Cartão};
        \draw[seta] (t) -- (s); \draw[seta] (s) -- (l); \draw[seta] (s) -- (b); \draw[seta] (l) -- (c);
        \draw[setaambar,dashed] (t.east) -| (b.north);
        \draw[setaambar,dashed] (t.west) -| ([xshift=-3mm]c.west) -- (c.west);
        \node[rotulo,text width=4cm,below=2mm of c.south east] {tracejado: componentes leem o tema direto do contexto, sem receber parâmetros};
      \end{tikzpicture}
    \end{column}
    \begin{column}{0.56\textwidth}
      \begin{definicao}[\emph{Tokens} de design]
        Decisões visuais nomeadas pelo \textbf{papel}, não pelo valor: \cod{primary}, \cod{surface}, \cod{titleMedium}. O componente pede ``a cor primária''; o tema decide qual é.
      \end{definicao}
      \begin{pordentro}[Contexto implícito]
        Tema, idioma e tamanho da tela ficam disponíveis para toda uma subárvore: \cod{InheritedWidget} no Flutter (\cod{Theme.of(context)}), \cod{CompositionLocal} no Compose. Quem lê é atualizado quando o valor muda.
      \end{pordentro}
      \begin{dica}
        Com \emph{tokens}, tema escuro e cor dinâmica (Android 12+) funcionam sem mexer nos componentes.
      \end{dica}
    \end{column}
  \end{columns}
\end{frame}

\begin{frame}{Pixels, densidade e telas de todos os tamanhos}
  \begin{columns}[T]
    \begin{column}{0.49\textwidth}
      \centering
      \begin{tikzpicture}[baseline=(current bounding box.north)]
        \draw[cinzamedio,step=0.2] (0,0) grid (1.2,1.2);
        \fill[ambar,opacity=0.7] (0.2,0.2) rectangle (1.0,1.0);
        \node[rotulo] at (0.6,-0.25) {160 dpi: 48\,dp = 48\,px};
        \draw[cinzamedio,step=0.0667] (2.6,0) grid (3.8,1.2);
        \fill[ambar,opacity=0.7] (2.8,0.2) rectangle (3.6,1.0);
        \node[rotulo] at (3.2,-0.25) {480 dpi: 48\,dp = 144\,px};
      \end{tikzpicture}
      \raggedright
      \begin{definicao}[Pixel independente de densidade]
        \cod{px = dp} $\times$ \cod{(dpi} $\div$ \cod{160)}: um botão de 48\,dp tem o mesmo tamanho físico em qualquer tela. No Flutter, é o pixel lógico (fator em \cod{devicePixelRatio}).
      \end{definicao}
      \begin{dica}[Texto]
        Fontes também seguem a escala escolhida pelo usuário: \cod{sp} no Android, \cod{TextScaler} no Flutter.
      \end{dica}
    \end{column}
    \begin{column}{0.49\textwidth}
      \begin{conceito}[Responsivo $\times$ adaptativo]
        \textbf{Responsivo}: o mesmo layout se estica ou reorganiza. \textbf{Adaptativo}: um layout por classe de tamanho de janela (barra inferior no celular, trilho lateral no tablet). Pense em \textbf{espaço disponível}, não em modelo de aparelho.
      \end{conceito}
      \begin{atencao}[Android 16 e telas grandes]
        Com \cod{targetSdk} 36 (exigido pela Play desde agosto de 2026), telas com largura mínima $\geq$ 600\,dp ignoram travas de orientação e de redimensionamento; no API 37 não haverá exceção.
      \end{atencao}
    \end{column}
  \end{columns}
\end{frame}

% =============================================================================
\section{Fundamentos: interfaces declarativas}

\begin{frame}{Interfaces imperativas e declarativas}
  \begin{columns}[T]
    \begin{column}{0.49\textwidth}
      \begin{definicao}[Interface imperativa]
        O programa cria os componentes uma vez e depois os \textbf{altera} passo a passo (\cod{setText}, \cod{setVisibility}). Basta esquecer um deles para a tela mostrar dados errados.
      \end{definicao}
    \end{column}
    \begin{column}{0.49\textwidth}
      \begin{definicao}[Interface declarativa]
        O programa \textbf{descreve} como a tela deve ser para cada estado. Quando o estado muda, o framework gera a nova descrição e aplica apenas as diferenças.
      \end{definicao}
    \end{column}
  \end{columns}
  \vspace{2pt}
  \centering
  \begin{tikzpicture}[node distance=2mm and 6mm,every node/.style={font=\scriptsize}]
    \node[rotulo,anchor=east,text width=1.6cm,align=right] (li) at (0,0) {\textbf{imperativo}};
    \node[caixa,right=3mm of li,minimum width=1.6cm] (e1) {estado muda};
    \node[caixaambar,right=of e1,minimum width=3.2cm] (c1) {seu código acha e altera\\ cada componente afetado};
    \node[caixa,right=of c1,minimum width=2.2cm] (t1) {tela (mutável)};
    \draw[seta] (e1) -- (c1); \draw[seta] (c1) -- (t1);
    \node[rotulo,anchor=east,text width=1.6cm,align=right] (ld) at (0,-1.05) {\textbf{declarativo}};
    \node[caixa,right=3mm of ld,minimum width=1.6cm] (e2) {estado muda};
    \node[caixapetroleo,right=of e2,minimum width=2.1cm] (f2) {UI = f(estado)};
    \node[caixa,right=of f2,minimum width=2.1cm] (d2) {nova descrição};
    \node[caixapetroleo,right=of d2,minimum width=2.6cm] (a2) {framework compara\\ e aplica diferenças};
    \draw[seta] (e2) -- (f2); \draw[seta] (f2) -- (d2); \draw[seta] (d2) -- (a2);
  \end{tikzpicture}
  \raggedright
  \begin{historia}
    Views do Android e UIKit do iOS (2008) são imperativos. O React (2013) popularizou a UI declarativa, depois adotada pelo Flutter (2018), SwiftUI (2019) e Jetpack Compose (2021).
  \end{historia}
\end{frame}

\begin{frame}{Composição: a tela como uma árvore}
  \begin{columns}[c]
    \begin{column}{0.52\textwidth}
      \centering
      \begin{tikzpicture}[level distance=10mm,
          every node/.style={caixa,font=\tiny,minimum height=5mm,inner sep=2pt},
          edge from parent/.style={draw=azulpetroleo,line width=0.6pt},
          level 1/.style={sibling distance=30mm},
          level 2/.style={sibling distance=21mm},
          level 3/.style={sibling distance=10mm}]
        \node[caixaescura,font=\tiny] (raiz) {TelaEquipamentos}
          child {node {BarraSuperior}}
          child {node {ListaEquipamentos}
            child {node[caixaambar,font=\tiny] {CartaoBancada}
              child {node {Ícone}}
              child {node {Textos}}
              child {node {Barra}}}
            child {node[caixaambar,font=\tiny] {CartaoBancada}}};
        \draw[seta] (-2.6,0.2) -- (-2.6,-3.2);
        \draw[setaambar] (-2.2,-3.2) -- (-2.2,0.2);
        \node[rotulo,anchor=north west,text width=6cm,align=left] at (-2.75,-3.45) {\textcolor{azulpetroleo}{$\downarrow$} dados descem por parâmetros\quad \textcolor{ambar}{$\uparrow$} eventos sobem por \emph{callbacks}};
      \end{tikzpicture}
    \end{column}
    \begin{column}{0.46\textwidth}
      \begin{definicao}[Composição]
        Construir interfaces complexas \textbf{combinando} componentes pequenos e independentes, em vez de herdar de componentes grandes. Cada componente recebe dados e devolve uma descrição de UI.
      \end{definicao}
      \begin{analogia}
        Como peças de montar: o cartão é feito de ícone, textos e indicador; a lista, de vários cartões. Trocar uma peça não obriga a refazer as outras.
      \end{analogia}
      \begin{dica}[Um bom componente]
        Faz uma coisa só, recebe tudo por parâmetro, não sabe de onde vêm os dados e pode ser visto isoladamente em uma \emph{preview}.
      \end{dica}
    \end{column}
  \end{columns}
\end{frame}

\begin{frame}{Layout por restrições}
  \begin{columns}[c]
    \begin{column}{0.47\textwidth}
      \centering
      \begin{tikzpicture}[node distance=11mm]
        \node[caixaescura,minimum width=3.4cm] (t) {Tela \quad 360 $\times$ 800};
        \node[caixa,minimum width=3.4cm,below=of t] (c) {Coluna};
        \node[caixaambar,minimum width=3.4cm,below=of c] (k) {Cartão};
        \draw[seta] ([xshift=-10mm]t.south) -- node[rotulo,left,text width=2.1cm,align=right]{largura 0--360\\ altura 0--800} ([xshift=-10mm]c.north);
        \draw[seta] ([xshift=-10mm]c.south) -- node[rotulo,left,text width=2.1cm,align=right]{largura 0--360\\ altura livre ($\infty$)} ([xshift=-10mm]k.north);
        \draw[setaambar] ([xshift=10mm]k.north) -- node[rotulo,right,text width=1.6cm,align=left]{``tenho\\ 360 $\times$ 96''} ([xshift=10mm]c.south);
        \node[rotulo,below=1.5mm of k,text width=4cm] {a coluna põe o cartão em $y = 0$\\ e o próximo em $y = 96$};
      \end{tikzpicture}
    \end{column}
    \begin{column}{0.51\textwidth}
      \begin{definicao}[Restrições (\emph{constraints})]
        Intervalo de tamanhos que o pai aceita para o filho: largura e altura mínimas e máximas. \textbf{Justa} quando mínimo = máximo; \textbf{frouxa} quando o mínimo é zero.
      \end{definicao}
      \begin{conceito}[O algoritmo em três frases]
        \textbf{Restrições descem}: o pai diz os limites. \textbf{Tamanhos sobem}: o filho escolhe um tamanho dentro deles. \textbf{O pai posiciona}: o filho não sabe onde está.
      \end{conceito}
      \begin{pordentro}[Uma passada só]
        Cada nó é medido uma única vez por quadro, o que mantém o layout rápido mesmo com milhares de nós. No Compose, cada \cod{Modifier} ajusta as restrições antes de repassá-las; no Flutter, \cod{Padding} e \cod{SizedBox} fazem esse papel.
      \end{pordentro}
    \end{column}
  \end{columns}
\end{frame}

\begin{frame}{Do estado aos pixels: o caminho de um quadro}
  \centering
  \begin{tikzpicture}[node distance=5mm,every node/.style={font=\scriptsize}]
    \node[caixaambar,minimum width=2cm,minimum height=11mm] (s) {estado\\ muda};
    \node[caixapetroleo,right=of s,minimum width=2.3cm,minimum height=11mm] (c) {\textbf{composição}\\ \tiny o que mostrar};
    \node[caixapetroleo,right=of c,minimum width=2.3cm,minimum height=11mm] (l) {\textbf{layout}\\ \tiny tamanho e posição};
    \node[caixapetroleo,right=of l,minimum width=2.3cm,minimum height=11mm] (d) {\textbf{desenho}\\ \tiny cores, textos, formas};
    \node[caixaescura,right=of d,minimum width=2.3cm,minimum height=11mm] (r) {\textbf{rasterização}\\ \tiny GPU gera os pixels};
    \draw[seta] (s) -- (c); \draw[seta] (c) -- (l); \draw[seta] (l) -- (d); \draw[seta] (d) -- (r);
    \node[rotulo,below=1mm of c] {Compose: \emph{Composition}\\ Flutter: \emph{build}};
    \node[rotulo,below=1mm of l] {\emph{Layout}\\ \emph{layout}};
    \node[rotulo,below=1mm of d] {\emph{Drawing}\\ \emph{paint} e \emph{composite}};
    \node[rotulo,below=1mm of r] {motor gráfico\\ (Skia, Impeller)};
  \end{tikzpicture}
  \raggedright
  \begin{columns}[T]
    \begin{column}{0.49\textwidth}
      \begin{definicao}[Quadro e orçamento de tempo]
        A tela é redesenhada no ritmo do sinal vertical do display (VSync): a 60\,Hz cada quadro tem \textbf{16,7\,ms}; a 120\,Hz, só \textbf{8,3\,ms}. Se o trabalho não couber, o quadro se perde e a animação engasga (\emph{jank}).
      \end{definicao}
    \end{column}
    \begin{column}{0.49\textwidth}
      \begin{pordentro}[Pular fases economiza tempo]
        Uma cor que anima só precisa de desenho; uma posição que muda, de layout e desenho. Ler o estado na fase mais tardia possível evita repetir as anteriores.
      \end{pordentro}
      \begin{dica}
        O \emph{Performance} do Flutter DevTools e o \emph{Layout Inspector} do Android Studio mostram onde o tempo do quadro foi gasto.
      \end{dica}
    \end{column}
  \end{columns}
\end{frame}

\begin{frame}{Por dentro: descrição, instância e renderização}
  \begin{columns}[T]
    \begin{column}{0.56\textwidth}
      \scriptsize
      \begin{tabularx}{\linewidth}{@{}L l l@{}}
        \toprule
        \cabecalho{Papel} & \cabecalho{Flutter} & \cabecalho{Compose} \\
        \midrule
        Descrição leve, recriada a cada mudança & \cod{Widget} & chamada \cod{@Composable} \\
        Instância duradoura, guarda o estado & \cod{Element} + \cod{State} & composição (\cod{remember}) \\
        Mede, posiciona e desenha & \cod{RenderObject} & \cod{LayoutNode} \\
        Informa leitores de tela & árvore de semântica & árvore de semântica \\
        \bottomrule
      \end{tabularx}
      \normalsize
      \begin{analogia}
        O widget é a planta do arquiteto: barata de redesenhar. O elemento é o apartamento construído, onde mora o estado. O objeto de renderização é a estrutura medida e pintada. Mudar a planta não obriga a demolir o prédio.
      \end{analogia}
    \end{column}
    \begin{column}{0.42\textwidth}
      \begin{conceito}[Por que separar]
        Descrições são criadas aos milhares por segundo, então precisam ser baratas. Medir e desenhar é caro, então os objetos de renderização são \textbf{reaproveitados} enquanto possível.
      \end{conceito}
      \begin{pordentro}[Semântica e acessibilidade]
        Em paralelo, o framework monta uma árvore de semântica (rótulos, papéis, ações) que o TalkBack e o VoiceOver leem. Ícone sem rótulo fica mudo para quem não enxerga (Módulo 04).
      \end{pordentro}
    \end{column}
  \end{columns}
\end{frame}

\begin{frame}{Atualizar sem refazer tudo}
  \begin{columns}[c]
    \begin{column}{0.42\textwidth}
      \centering
      \begin{tikzpicture}[level distance=10mm,
          every node/.style={caixa,font=\tiny,minimum height=5mm,minimum width=9mm,inner sep=2pt},
          edge from parent/.style={draw=azulpetroleo,line width=0.6pt},
          level 1/.style={sibling distance=22mm},
          level 2/.style={sibling distance=11mm}]
        \node {Tela}
          child {node {Barra} child {node {Título}}}
          child {node[caixaambarforte] (m) {Medidor}
            child {node[caixaambar] {Valor}}
            child {node[caixaambar] {Barra}}};
        \node[rotulo,below=11mm of m,text width=3.4cm] {o estado \cod{temperatura} mudou: só o \textbf{Medidor} e seus filhos são reconstruídos};
      \end{tikzpicture}
    \end{column}
    \begin{column}{0.56\textwidth}
      \begin{conceito}[Flutter: \emph{rebuild} e reconciliação]
        \cod{setState} marca o elemento como sujo; no próximo quadro, só ele executa \cod{build}. Filhos de mesmo tipo e chave são reaproveitados e só atualizados.
      \end{conceito}
      \begin{conceito}[Compose: recomposição inteligente]
        O Compose sabe quais funções leram cada estado e só executa essas de novo; as que recebem parâmetros iguais aos anteriores são puladas (\emph{skipping}).
      \end{conceito}
      \begin{atencao}[Consequência importante]
        Funções de interface podem rodar muitas vezes, em qualquer ordem: devem ser \textbf{rápidas e sem efeitos colaterais}. Rede e banco ficam em APIs de efeito (Módulo 07).
      \end{atencao}
    \end{column}
  \end{columns}
\end{frame}

\begin{frame}[fragile]{Identidade e chaves}
  \begin{columns}[T]
    \begin{column}{0.47\textwidth}
      \centering
      \begin{tikzpicture}[baseline=(current bounding box.north),every node/.style={font=\tiny}]
        \tikzset{item/.style={caixa,minimum width=1.5cm,minimum height=5mm,inner sep=1pt}}
        \node[rotulo] at (0.75,0.55) {antes};
        \node[item,caixaambar] (a) at (0.75,0) {A \quad expandido};
        \node[item] at (0.75,-0.6) {B};
        \node[item] at (0.75,-1.2) {C};
        \node[rotulo,text width=2cm] at (2.95,0.55) {remove A, \textbf{sem} chave};
        \node[item,caixaambar,draw=vermelhoerro] at (2.95,0) {B \quad expandido};
        \node[item] at (2.95,-0.6) {C};
        \node[rotulo,text width=2cm] at (5.15,0.55) {remove A, \textbf{com} chave};
        \node[item] at (5.15,0) {B};
        \node[item] at (5.15,-0.6) {C};
        \node[rotulo,text width=2cm,text=vermelhoerro] at (2.95,-1.25) {o estado ficou na posição 1};
        \node[rotulo,text width=2cm,text=verdeok] at (5.15,-1.25) {o estado saiu junto com A};
      \end{tikzpicture}
\begin{codigo}[title=Compose]{Kotlin}
items(itens, key = { it.id }) { eq ->
  ItemEquipamento(eq)
}
\end{codigo}
\begin{codigo}[title=Flutter]{Dart}
ItemEquipamento(
    key: ValueKey(eq.id), eq: eq, onTap: abrir)
\end{codigo}
    \end{column}
    \begin{column}{0.51\textwidth}
      \begin{definicao}[Chave (\emph{key})]
        Identificador estável que diz ao framework qual item é qual entre duas versões da árvore. Sem chave, a correspondência é feita pela \textbf{posição}.
      \end{definicao}
      \begin{conceito}[Quando importa]
        Em listas cujos itens têm estado próprio (expandido, campo digitado, animação) e podem ser inseridos, removidos ou reordenados.
      \end{conceito}
      \begin{dica}
        Use um identificador do domínio, como o \cod{id} do equipamento. Nunca o índice da lista nem um valor aleatório criado a cada reconstrução.
      \end{dica}
    \end{column}
  \end{columns}
\end{frame}

\begin{frame}{Temas, \emph{tokens} e contexto implícito}
  \begin{columns}[c]
    \begin{column}{0.42\textwidth}
      \centering
      \begin{tikzpicture}[node distance=7mm and 3mm,every node/.style={font=\tiny}]
        \node[caixaambarforte,minimum width=2.4cm,font=\tiny\bfseries] (t) {Tema};
        \node[caixa,below=of t,minimum width=1.6cm,font=\tiny] (s) {Tela};
        \node[caixa,below left=of s,minimum width=1.3cm,font=\tiny] (l) {Lista};
        \node[caixa,below right=of s,minimum width=1.3cm,font=\tiny] (b) {Botão};
        \node[caixa,below=of l,minimum width=1.3cm,font=\tiny] (c) {Cartão};
        \draw[seta] (t) -- (s); \draw[seta] (s) -- (l); \draw[seta] (s) -- (b); \draw[seta] (l) -- (c);
        \draw[setaambar,dashed] (t.east) -| (b.north);
        \draw[setaambar,dashed] (t.west) -| ([xshift=-3mm]c.west) -- (c.west);
        \node[rotulo,text width=4cm,below=2mm of c.south east] {tracejado: componentes leem o tema direto do contexto, sem receber parâmetros};
      \end{tikzpicture}
    \end{column}
    \begin{column}{0.56\textwidth}
      \begin{definicao}[\emph{Tokens} de design]
        Decisões visuais nomeadas pelo \textbf{papel}, não pelo valor: \cod{primary}, \cod{surface}, \cod{titleMedium}. O componente pede ``a cor primária''; o tema decide qual é.
      \end{definicao}
      \begin{pordentro}[Contexto implícito]
        Tema, idioma e tamanho da tela ficam disponíveis para toda uma subárvore: \cod{InheritedWidget} no Flutter (\cod{Theme.of(context)}), \cod{CompositionLocal} no Compose. Quem lê é atualizado quando o valor muda.
      \end{pordentro}
      \begin{dica}
        Com \emph{tokens}, tema escuro e cor dinâmica (Android 12+) funcionam sem mexer nos componentes.
      \end{dica}
    \end{column}
  \end{columns}
\end{frame}

\begin{frame}{Pixels, densidade e telas de todos os tamanhos}
  \begin{columns}[T]
    \begin{column}{0.49\textwidth}
      \centering
      \begin{tikzpicture}[baseline=(current bounding box.north)]
        \draw[cinzamedio,step=0.2] (0,0) grid (1.2,1.2);
        \fill[ambar,opacity=0.7] (0.2,0.2) rectangle (1.0,1.0);
        \node[rotulo] at (0.6,-0.25) {160 dpi: 48\,dp = 48\,px};
        \draw[cinzamedio,step=0.0667] (2.6,0) grid (3.8,1.2);
        \fill[ambar,opacity=0.7] (2.8,0.2) rectangle (3.6,1.0);
        \node[rotulo] at (3.2,-0.25) {480 dpi: 48\,dp = 144\,px};
      \end{tikzpicture}
      \raggedright
      \begin{definicao}[Pixel independente de densidade]
        \cod{px = dp} $\times$ \cod{(dpi} $\div$ \cod{160)}: um botão de 48\,dp tem o mesmo tamanho físico em qualquer tela. No Flutter, é o pixel lógico (fator em \cod{devicePixelRatio}).
      \end{definicao}
      \begin{dica}[Texto]
        Fontes também seguem a escala escolhida pelo usuário: \cod{sp} no Android, \cod{TextScaler} no Flutter.
      \end{dica}
    \end{column}
    \begin{column}{0.49\textwidth}
      \begin{conceito}[Responsivo $\times$ adaptativo]
        \textbf{Responsivo}: o mesmo layout se estica ou reorganiza. \textbf{Adaptativo}: um layout por classe de tamanho de janela (barra inferior no celular, trilho lateral no tablet). Pense em \textbf{espaço disponível}, não em modelo de aparelho.
      \end{conceito}
      \begin{atencao}[Android 16 e telas grandes]
        Com \cod{targetSdk} 36 (exigido pela Play desde agosto de 2026), telas com largura mínima $\geq$ 600\,dp ignoram travas de orientação e de redimensionamento; no API 37 não haverá exceção.
      \end{atencao}
    \end{column}
  \end{columns}
\end{frame}

% =============================================================================
\section{Fundamentos: interfaces declarativas}

\begin{frame}{Interfaces imperativas e declarativas}
  \begin{columns}[T]
    \begin{column}{0.49\textwidth}
      \begin{definicao}[Interface imperativa]
        O programa cria os componentes uma vez e depois os \textbf{altera} passo a passo (\cod{setText}, \cod{setVisibility}). Basta esquecer um deles para a tela mostrar dados errados.
      \end{definicao}
    \end{column}
    \begin{column}{0.49\textwidth}
      \begin{definicao}[Interface declarativa]
        O programa \textbf{descreve} como a tela deve ser para cada estado. Quando o estado muda, o framework gera a nova descrição e aplica apenas as diferenças.
      \end{definicao}
    \end{column}
  \end{columns}
  \vspace{2pt}
  \centering
  \begin{tikzpicture}[node distance=2mm and 6mm,every node/.style={font=\scriptsize}]
    \node[rotulo,anchor=east,text width=1.6cm,align=right] (li) at (0,0) {\textbf{imperativo}};
    \node[caixa,right=3mm of li,minimum width=1.6cm] (e1) {estado muda};
    \node[caixaambar,right=of e1,minimum width=3.2cm] (c1) {seu código acha e altera\\ cada componente afetado};
    \node[caixa,right=of c1,minimum width=2.2cm] (t1) {tela (mutável)};
    \draw[seta] (e1) -- (c1); \draw[seta] (c1) -- (t1);
    \node[rotulo,anchor=east,text width=1.6cm,align=right] (ld) at (0,-1.05) {\textbf{declarativo}};
    \node[caixa,right=3mm of ld,minimum width=1.6cm] (e2) {estado muda};
    \node[caixapetroleo,right=of e2,minimum width=2.1cm] (f2) {UI = f(estado)};
    \node[caixa,right=of f2,minimum width=2.1cm] (d2) {nova descrição};
    \node[caixapetroleo,right=of d2,minimum width=2.6cm] (a2) {framework compara\\ e aplica diferenças};
    \draw[seta] (e2) -- (f2); \draw[seta] (f2) -- (d2); \draw[seta] (d2) -- (a2);
  \end{tikzpicture}
  \raggedright
  \begin{historia}
    Views do Android e UIKit do iOS (2008) são imperativos. O React (2013) popularizou a UI declarativa, depois adotada pelo Flutter (2018), SwiftUI (2019) e Jetpack Compose (2021).
  \end{historia}
\end{frame}

\begin{frame}{Composição: a tela como uma árvore}
  \begin{columns}[c]
    \begin{column}{0.52\textwidth}
      \centering
      \begin{tikzpicture}[level distance=10mm,
          every node/.style={caixa,font=\tiny,minimum height=5mm,inner sep=2pt},
          edge from parent/.style={draw=azulpetroleo,line width=0.6pt},
          level 1/.style={sibling distance=30mm},
          level 2/.style={sibling distance=21mm},
          level 3/.style={sibling distance=10mm}]
        \node[caixaescura,font=\tiny] (raiz) {TelaEquipamentos}
          child {node {BarraSuperior}}
          child {node {ListaEquipamentos}
            child {node[caixaambar,font=\tiny] {CartaoBancada}
              child {node {Ícone}}
              child {node {Textos}}
              child {node {Barra}}}
            child {node[caixaambar,font=\tiny] {CartaoBancada}}};
        \draw[seta] (-2.6,0.2) -- (-2.6,-3.2);
        \draw[setaambar] (-2.2,-3.2) -- (-2.2,0.2);
        \node[rotulo,anchor=north west,text width=6cm,align=left] at (-2.75,-3.45) {\textcolor{azulpetroleo}{$\downarrow$} dados descem por parâmetros\quad \textcolor{ambar}{$\uparrow$} eventos sobem por \emph{callbacks}};
      \end{tikzpicture}
    \end{column}
    \begin{column}{0.46\textwidth}
      \begin{definicao}[Composição]
        Construir interfaces complexas \textbf{combinando} componentes pequenos e independentes, em vez de herdar de componentes grandes. Cada componente recebe dados e devolve uma descrição de UI.
      \end{definicao}
      \begin{analogia}
        Como peças de montar: o cartão é feito de ícone, textos e indicador; a lista, de vários cartões. Trocar uma peça não obriga a refazer as outras.
      \end{analogia}
      \begin{dica}[Um bom componente]
        Faz uma coisa só, recebe tudo por parâmetro, não sabe de onde vêm os dados e pode ser visto isoladamente em uma \emph{preview}.
      \end{dica}
    \end{column}
  \end{columns}
\end{frame}

\begin{frame}{Layout por restrições}
  \begin{columns}[c]
    \begin{column}{0.47\textwidth}
      \centering
      \begin{tikzpicture}[node distance=11mm]
        \node[caixaescura,minimum width=3.4cm] (t) {Tela \quad 360 $\times$ 800};
        \node[caixa,minimum width=3.4cm,below=of t] (c) {Coluna};
        \node[caixaambar,minimum width=3.4cm,below=of c] (k) {Cartão};
        \draw[seta] ([xshift=-10mm]t.south) -- node[rotulo,left,text width=2.1cm,align=right]{largura 0--360\\ altura 0--800} ([xshift=-10mm]c.north);
        \draw[seta] ([xshift=-10mm]c.south) -- node[rotulo,left,text width=2.1cm,align=right]{largura 0--360\\ altura livre ($\infty$)} ([xshift=-10mm]k.north);
        \draw[setaambar] ([xshift=10mm]k.north) -- node[rotulo,right,text width=1.6cm,align=left]{``tenho\\ 360 $\times$ 96''} ([xshift=10mm]c.south);
        \node[rotulo,below=1.5mm of k,text width=4cm] {a coluna põe o cartão em $y = 0$\\ e o próximo em $y = 96$};
      \end{tikzpicture}
    \end{column}
    \begin{column}{0.51\textwidth}
      \begin{definicao}[Restrições (\emph{constraints})]
        Intervalo de tamanhos que o pai aceita para o filho: largura e altura mínimas e máximas. \textbf{Justa} quando mínimo = máximo; \textbf{frouxa} quando o mínimo é zero.
      \end{definicao}
      \begin{conceito}[O algoritmo em três frases]
        \textbf{Restrições descem}: o pai diz os limites. \textbf{Tamanhos sobem}: o filho escolhe um tamanho dentro deles. \textbf{O pai posiciona}: o filho não sabe onde está.
      \end{conceito}
      \begin{pordentro}[Uma passada só]
        Cada nó é medido uma única vez por quadro, o que mantém o layout rápido mesmo com milhares de nós. No Compose, cada \cod{Modifier} ajusta as restrições antes de repassá-las; no Flutter, \cod{Padding} e \cod{SizedBox} fazem esse papel.
      \end{pordentro}
    \end{column}
  \end{columns}
\end{frame}

\begin{frame}{Do estado aos pixels: o caminho de um quadro}
  \centering
  \begin{tikzpicture}[node distance=5mm,every node/.style={font=\scriptsize}]
    \node[caixaambar,minimum width=2cm,minimum height=11mm] (s) {estado\\ muda};
    \node[caixapetroleo,right=of s,minimum width=2.3cm,minimum height=11mm] (c) {\textbf{composição}\\ \tiny o que mostrar};
    \node[caixapetroleo,right=of c,minimum width=2.3cm,minimum height=11mm] (l) {\textbf{layout}\\ \tiny tamanho e posição};
    \node[caixapetroleo,right=of l,minimum width=2.3cm,minimum height=11mm] (d) {\textbf{desenho}\\ \tiny cores, textos, formas};
    \node[caixaescura,right=of d,minimum width=2.3cm,minimum height=11mm] (r) {\textbf{rasterização}\\ \tiny GPU gera os pixels};
    \draw[seta] (s) -- (c); \draw[seta] (c) -- (l); \draw[seta] (l) -- (d); \draw[seta] (d) -- (r);
    \node[rotulo,below=1mm of c] {Compose: \emph{Composition}\\ Flutter: \emph{build}};
    \node[rotulo,below=1mm of l] {\emph{Layout}\\ \emph{layout}};
    \node[rotulo,below=1mm of d] {\emph{Drawing}\\ \emph{paint} e \emph{composite}};
    \node[rotulo,below=1mm of r] {motor gráfico\\ (Skia, Impeller)};
  \end{tikzpicture}
  \raggedright
  \begin{columns}[T]
    \begin{column}{0.49\textwidth}
      \begin{definicao}[Quadro e orçamento de tempo]
        A tela é redesenhada no ritmo do sinal vertical do display (VSync): a 60\,Hz cada quadro tem \textbf{16,7\,ms}; a 120\,Hz, só \textbf{8,3\,ms}. Se o trabalho não couber, o quadro se perde e a animação engasga (\emph{jank}).
      \end{definicao}
    \end{column}
    \begin{column}{0.49\textwidth}
      \begin{pordentro}[Pular fases economiza tempo]
        Uma cor que anima só precisa de desenho; uma posição que muda, de layout e desenho. Ler o estado na fase mais tardia possível evita repetir as anteriores.
      \end{pordentro}
      \begin{dica}
        O \emph{Performance} do Flutter DevTools e o \emph{Layout Inspector} do Android Studio mostram onde o tempo do quadro foi gasto.
      \end{dica}
    \end{column}
  \end{columns}
\end{frame}

\begin{frame}{Por dentro: descrição, instância e renderização}
  \begin{columns}[T]
    \begin{column}{0.56\textwidth}
      \scriptsize
      \begin{tabularx}{\linewidth}{@{}L l l@{}}
        \toprule
        \cabecalho{Papel} & \cabecalho{Flutter} & \cabecalho{Compose} \\
        \midrule
        Descrição leve, recriada a cada mudança & \cod{Widget} & chamada \cod{@Composable} \\
        Instância duradoura, guarda o estado & \cod{Element} + \cod{State} & composição (\cod{remember}) \\
        Mede, posiciona e desenha & \cod{RenderObject} & \cod{LayoutNode} \\
        Informa leitores de tela & árvore de semântica & árvore de semântica \\
        \bottomrule
      \end{tabularx}
      \normalsize
      \begin{analogia}
        O widget é a planta do arquiteto: barata de redesenhar. O elemento é o apartamento construído, onde mora o estado. O objeto de renderização é a estrutura medida e pintada. Mudar a planta não obriga a demolir o prédio.
      \end{analogia}
    \end{column}
    \begin{column}{0.42\textwidth}
      \begin{conceito}[Por que separar]
        Descrições são criadas aos milhares por segundo, então precisam ser baratas. Medir e desenhar é caro, então os objetos de renderização são \textbf{reaproveitados} enquanto possível.
      \end{conceito}
      \begin{pordentro}[Semântica e acessibilidade]
        Em paralelo, o framework monta uma árvore de semântica (rótulos, papéis, ações) que o TalkBack e o VoiceOver leem. Ícone sem rótulo fica mudo para quem não enxerga (Módulo 04).
      \end{pordentro}
    \end{column}
  \end{columns}
\end{frame}

\begin{frame}{Atualizar sem refazer tudo}
  \begin{columns}[c]
    \begin{column}{0.42\textwidth}
      \centering
      \begin{tikzpicture}[level distance=10mm,
          every node/.style={caixa,font=\tiny,minimum height=5mm,minimum width=9mm,inner sep=2pt},
          edge from parent/.style={draw=azulpetroleo,line width=0.6pt},
          level 1/.style={sibling distance=22mm},
          level 2/.style={sibling distance=11mm}]
        \node {Tela}
          child {node {Barra} child {node {Título}}}
          child {node[caixaambarforte] (m) {Medidor}
            child {node[caixaambar] {Valor}}
            child {node[caixaambar] {Barra}}};
        \node[rotulo,below=11mm of m,text width=3.4cm] {o estado \cod{temperatura} mudou: só o \textbf{Medidor} e seus filhos são reconstruídos};
      \end{tikzpicture}
    \end{column}
    \begin{column}{0.56\textwidth}
      \begin{conceito}[Flutter: \emph{rebuild} e reconciliação]
        \cod{setState} marca o elemento como sujo; no próximo quadro, só ele executa \cod{build}. Filhos de mesmo tipo e chave são reaproveitados e só atualizados.
      \end{conceito}
      \begin{conceito}[Compose: recomposição inteligente]
        O Compose sabe quais funções leram cada estado e só executa essas de novo; as que recebem parâmetros iguais aos anteriores são puladas (\emph{skipping}).
      \end{conceito}
      \begin{atencao}[Consequência importante]
        Funções de interface podem rodar muitas vezes, em qualquer ordem: devem ser \textbf{rápidas e sem efeitos colaterais}. Rede e banco ficam em APIs de efeito (Módulo 07).
      \end{atencao}
    \end{column}
  \end{columns}
\end{frame}

\begin{frame}[fragile]{Identidade e chaves}
  \begin{columns}[T]
    \begin{column}{0.47\textwidth}
      \centering
      \begin{tikzpicture}[baseline=(current bounding box.north),every node/.style={font=\tiny}]
        \tikzset{item/.style={caixa,minimum width=1.5cm,minimum height=5mm,inner sep=1pt}}
        \node[rotulo] at (0.75,0.55) {antes};
        \node[item,caixaambar] (a) at (0.75,0) {A \quad expandido};
        \node[item] at (0.75,-0.6) {B};
        \node[item] at (0.75,-1.2) {C};
        \node[rotulo,text width=2cm] at (2.95,0.55) {remove A, \textbf{sem} chave};
        \node[item,caixaambar,draw=vermelhoerro] at (2.95,0) {B \quad expandido};
        \node[item] at (2.95,-0.6) {C};
        \node[rotulo,text width=2cm] at (5.15,0.55) {remove A, \textbf{com} chave};
        \node[item] at (5.15,0) {B};
        \node[item] at (5.15,-0.6) {C};
        \node[rotulo,text width=2cm,text=vermelhoerro] at (2.95,-1.25) {o estado ficou na posição 1};
        \node[rotulo,text width=2cm,text=verdeok] at (5.15,-1.25) {o estado saiu junto com A};
      \end{tikzpicture}
\begin{codigo}[title=Compose]{Kotlin}
items(itens, key = { it.id }) { eq ->
  ItemEquipamento(eq)
}
\end{codigo}
\begin{codigo}[title=Flutter]{Dart}
ItemEquipamento(
    key: ValueKey(eq.id), eq: eq, onTap: abrir)
\end{codigo}
    \end{column}
    \begin{column}{0.51\textwidth}
      \begin{definicao}[Chave (\emph{key})]
        Identificador estável que diz ao framework qual item é qual entre duas versões da árvore. Sem chave, a correspondência é feita pela \textbf{posição}.
      \end{definicao}
      \begin{conceito}[Quando importa]
        Em listas cujos itens têm estado próprio (expandido, campo digitado, animação) e podem ser inseridos, removidos ou reordenados.
      \end{conceito}
      \begin{dica}
        Use um identificador do domínio, como o \cod{id} do equipamento. Nunca o índice da lista nem um valor aleatório criado a cada reconstrução.
      \end{dica}
    \end{column}
  \end{columns}
\end{frame}

\begin{frame}{Temas, \emph{tokens} e contexto implícito}
  \begin{columns}[c]
    \begin{column}{0.42\textwidth}
      \centering
      \begin{tikzpicture}[node distance=7mm and 3mm,every node/.style={font=\tiny}]
        \node[caixaambarforte,minimum width=2.4cm,font=\tiny\bfseries] (t) {Tema};
        \node[caixa,below=of t,minimum width=1.6cm,font=\tiny] (s) {Tela};
        \node[caixa,below left=of s,minimum width=1.3cm,font=\tiny] (l) {Lista};
        \node[caixa,below right=of s,minimum width=1.3cm,font=\tiny] (b) {Botão};
        \node[caixa,below=of l,minimum width=1.3cm,font=\tiny] (c) {Cartão};
        \draw[seta] (t) -- (s); \draw[seta] (s) -- (l); \draw[seta] (s) -- (b); \draw[seta] (l) -- (c);
        \draw[setaambar,dashed] (t.east) -| (b.north);
        \draw[setaambar,dashed] (t.west) -| ([xshift=-3mm]c.west) -- (c.west);
        \node[rotulo,text width=4cm,below=2mm of c.south east] {tracejado: componentes leem o tema direto do contexto, sem receber parâmetros};
      \end{tikzpicture}
    \end{column}
    \begin{column}{0.56\textwidth}
      \begin{definicao}[\emph{Tokens} de design]
        Decisões visuais nomeadas pelo \textbf{papel}, não pelo valor: \cod{primary}, \cod{surface}, \cod{titleMedium}. O componente pede ``a cor primária''; o tema decide qual é.
      \end{definicao}
      \begin{pordentro}[Contexto implícito]
        Tema, idioma e tamanho da tela ficam disponíveis para toda uma subárvore: \cod{InheritedWidget} no Flutter (\cod{Theme.of(context)}), \cod{CompositionLocal} no Compose. Quem lê é atualizado quando o valor muda.
      \end{pordentro}
      \begin{dica}
        Com \emph{tokens}, tema escuro e cor dinâmica (Android 12+) funcionam sem mexer nos componentes.
      \end{dica}
    \end{column}
  \end{columns}
\end{frame}

\begin{frame}{Pixels, densidade e telas de todos os tamanhos}
  \begin{columns}[T]
    \begin{column}{0.49\textwidth}
      \centering
      \begin{tikzpicture}[baseline=(current bounding box.north)]
        \draw[cinzamedio,step=0.2] (0,0) grid (1.2,1.2);
        \fill[ambar,opacity=0.7] (0.2,0.2) rectangle (1.0,1.0);
        \node[rotulo] at (0.6,-0.25) {160 dpi: 48\,dp = 48\,px};
        \draw[cinzamedio,step=0.0667] (2.6,0) grid (3.8,1.2);
        \fill[ambar,opacity=0.7] (2.8,0.2) rectangle (3.6,1.0);
        \node[rotulo] at (3.2,-0.25) {480 dpi: 48\,dp = 144\,px};
      \end{tikzpicture}
      \raggedright
      \begin{definicao}[Pixel independente de densidade]
        \cod{px = dp} $\times$ \cod{(dpi} $\div$ \cod{160)}: um botão de 48\,dp tem o mesmo tamanho físico em qualquer tela. No Flutter, é o pixel lógico (fator em \cod{devicePixelRatio}).
      \end{definicao}
      \begin{dica}[Texto]
        Fontes também seguem a escala escolhida pelo usuário: \cod{sp} no Android, \cod{TextScaler} no Flutter.
      \end{dica}
    \end{column}
    \begin{column}{0.49\textwidth}
      \begin{conceito}[Responsivo $\times$ adaptativo]
        \textbf{Responsivo}: o mesmo layout se estica ou reorganiza. \textbf{Adaptativo}: um layout por classe de tamanho de janela (barra inferior no celular, trilho lateral no tablet). Pense em \textbf{espaço disponível}, não em modelo de aparelho.
      \end{conceito}
      \begin{atencao}[Android 16 e telas grandes]
        Com \cod{targetSdk} 36 (exigido pela Play desde agosto de 2026), telas com largura mínima $\geq$ 600\,dp ignoram travas de orientação e de redimensionamento; no API 37 não haverá exceção.
      \end{atencao}
    \end{column}
  \end{columns}
\end{frame}

% =============================================================================
\section{Fundamentos: interfaces declarativas}

\begin{frame}{Interfaces imperativas e declarativas}
  \begin{columns}[T]
    \begin{column}{0.49\textwidth}
      \begin{definicao}[Interface imperativa]
        O programa cria os componentes uma vez e depois os \textbf{altera} passo a passo (\cod{setText}, \cod{setVisibility}). Basta esquecer um deles para a tela mostrar dados errados.
      \end{definicao}
    \end{column}
    \begin{column}{0.49\textwidth}
      \begin{definicao}[Interface declarativa]
        O programa \textbf{descreve} como a tela deve ser para cada estado. Quando o estado muda, o framework gera a nova descrição e aplica apenas as diferenças.
      \end{definicao}
    \end{column}
  \end{columns}
  \vspace{2pt}
  \centering
  \begin{tikzpicture}[node distance=2mm and 6mm,every node/.style={font=\scriptsize}]
    \node[rotulo,anchor=east,text width=1.6cm,align=right] (li) at (0,0) {\textbf{imperativo}};
    \node[caixa,right=3mm of li,minimum width=1.6cm] (e1) {estado muda};
    \node[caixaambar,right=of e1,minimum width=3.2cm] (c1) {seu código acha e altera\\ cada componente afetado};
    \node[caixa,right=of c1,minimum width=2.2cm] (t1) {tela (mutável)};
    \draw[seta] (e1) -- (c1); \draw[seta] (c1) -- (t1);
    \node[rotulo,anchor=east,text width=1.6cm,align=right] (ld) at (0,-1.05) {\textbf{declarativo}};
    \node[caixa,right=3mm of ld,minimum width=1.6cm] (e2) {estado muda};
    \node[caixapetroleo,right=of e2,minimum width=2.1cm] (f2) {UI = f(estado)};
    \node[caixa,right=of f2,minimum width=2.1cm] (d2) {nova descrição};
    \node[caixapetroleo,right=of d2,minimum width=2.6cm] (a2) {framework compara\\ e aplica diferenças};
    \draw[seta] (e2) -- (f2); \draw[seta] (f2) -- (d2); \draw[seta] (d2) -- (a2);
  \end{tikzpicture}
  \raggedright
  \begin{historia}
    Views do Android e UIKit do iOS (2008) são imperativos. O React (2013) popularizou a UI declarativa, depois adotada pelo Flutter (2018), SwiftUI (2019) e Jetpack Compose (2021).
  \end{historia}
\end{frame}

\begin{frame}{Composição: a tela como uma árvore}
  \begin{columns}[c]
    \begin{column}{0.52\textwidth}
      \centering
      \begin{tikzpicture}[level distance=10mm,
          every node/.style={caixa,font=\tiny,minimum height=5mm,inner sep=2pt},
          edge from parent/.style={draw=azulpetroleo,line width=0.6pt},
          level 1/.style={sibling distance=30mm},
          level 2/.style={sibling distance=21mm},
          level 3/.style={sibling distance=10mm}]
        \node[caixaescura,font=\tiny] (raiz) {TelaEquipamentos}
          child {node {BarraSuperior}}
          child {node {ListaEquipamentos}
            child {node[caixaambar,font=\tiny] {CartaoBancada}
              child {node {Ícone}}
              child {node {Textos}}
              child {node {Barra}}}
            child {node[caixaambar,font=\tiny] {CartaoBancada}}};
        \draw[seta] (-2.6,0.2) -- (-2.6,-3.2);
        \draw[setaambar] (-2.2,-3.2) -- (-2.2,0.2);
        \node[rotulo,anchor=north west,text width=6cm,align=left] at (-2.75,-3.45) {\textcolor{azulpetroleo}{$\downarrow$} dados descem por parâmetros\quad \textcolor{ambar}{$\uparrow$} eventos sobem por \emph{callbacks}};
      \end{tikzpicture}
    \end{column}
    \begin{column}{0.46\textwidth}
      \begin{definicao}[Composição]
        Construir interfaces complexas \textbf{combinando} componentes pequenos e independentes, em vez de herdar de componentes grandes. Cada componente recebe dados e devolve uma descrição de UI.
      \end{definicao}
      \begin{analogia}
        Como peças de montar: o cartão é feito de ícone, textos e indicador; a lista, de vários cartões. Trocar uma peça não obriga a refazer as outras.
      \end{analogia}
      \begin{dica}[Um bom componente]
        Faz uma coisa só, recebe tudo por parâmetro, não sabe de onde vêm os dados e pode ser visto isoladamente em uma \emph{preview}.
      \end{dica}
    \end{column}
  \end{columns}
\end{frame}

\begin{frame}{Layout por restrições}
  \begin{columns}[c]
    \begin{column}{0.47\textwidth}
      \centering
      \begin{tikzpicture}[node distance=11mm]
        \node[caixaescura,minimum width=3.4cm] (t) {Tela \quad 360 $\times$ 800};
        \node[caixa,minimum width=3.4cm,below=of t] (c) {Coluna};
        \node[caixaambar,minimum width=3.4cm,below=of c] (k) {Cartão};
        \draw[seta] ([xshift=-10mm]t.south) -- node[rotulo,left,text width=2.1cm,align=right]{largura 0--360\\ altura 0--800} ([xshift=-10mm]c.north);
        \draw[seta] ([xshift=-10mm]c.south) -- node[rotulo,left,text width=2.1cm,align=right]{largura 0--360\\ altura livre ($\infty$)} ([xshift=-10mm]k.north);
        \draw[setaambar] ([xshift=10mm]k.north) -- node[rotulo,right,text width=1.6cm,align=left]{``tenho\\ 360 $\times$ 96''} ([xshift=10mm]c.south);
        \node[rotulo,below=1.5mm of k,text width=4cm] {a coluna põe o cartão em $y = 0$\\ e o próximo em $y = 96$};
      \end{tikzpicture}
    \end{column}
    \begin{column}{0.51\textwidth}
      \begin{definicao}[Restrições (\emph{constraints})]
        Intervalo de tamanhos que o pai aceita para o filho: largura e altura mínimas e máximas. \textbf{Justa} quando mínimo = máximo; \textbf{frouxa} quando o mínimo é zero.
      \end{definicao}
      \begin{conceito}[O algoritmo em três frases]
        \textbf{Restrições descem}: o pai diz os limites. \textbf{Tamanhos sobem}: o filho escolhe um tamanho dentro deles. \textbf{O pai posiciona}: o filho não sabe onde está.
      \end{conceito}
      \begin{pordentro}[Uma passada só]
        Cada nó é medido uma única vez por quadro, o que mantém o layout rápido mesmo com milhares de nós. No Compose, cada \cod{Modifier} ajusta as restrições antes de repassá-las; no Flutter, \cod{Padding} e \cod{SizedBox} fazem esse papel.
      \end{pordentro}
    \end{column}
  \end{columns}
\end{frame}

\begin{frame}{Do estado aos pixels: o caminho de um quadro}
  \centering
  \begin{tikzpicture}[node distance=5mm,every node/.style={font=\scriptsize}]
    \node[caixaambar,minimum width=2cm,minimum height=11mm] (s) {estado\\ muda};
    \node[caixapetroleo,right=of s,minimum width=2.3cm,minimum height=11mm] (c) {\textbf{composição}\\ \tiny o que mostrar};
    \node[caixapetroleo,right=of c,minimum width=2.3cm,minimum height=11mm] (l) {\textbf{layout}\\ \tiny tamanho e posição};
    \node[caixapetroleo,right=of l,minimum width=2.3cm,minimum height=11mm] (d) {\textbf{desenho}\\ \tiny cores, textos, formas};
    \node[caixaescura,right=of d,minimum width=2.3cm,minimum height=11mm] (r) {\textbf{rasterização}\\ \tiny GPU gera os pixels};
    \draw[seta] (s) -- (c); \draw[seta] (c) -- (l); \draw[seta] (l) -- (d); \draw[seta] (d) -- (r);
    \node[rotulo,below=1mm of c] {Compose: \emph{Composition}\\ Flutter: \emph{build}};
    \node[rotulo,below=1mm of l] {\emph{Layout}\\ \emph{layout}};
    \node[rotulo,below=1mm of d] {\emph{Drawing}\\ \emph{paint} e \emph{composite}};
    \node[rotulo,below=1mm of r] {motor gráfico\\ (Skia, Impeller)};
  \end{tikzpicture}
  \raggedright
  \begin{columns}[T]
    \begin{column}{0.49\textwidth}
      \begin{definicao}[Quadro e orçamento de tempo]
        A tela é redesenhada no ritmo do sinal vertical do display (VSync): a 60\,Hz cada quadro tem \textbf{16,7\,ms}; a 120\,Hz, só \textbf{8,3\,ms}. Se o trabalho não couber, o quadro se perde e a animação engasga (\emph{jank}).
      \end{definicao}
    \end{column}
    \begin{column}{0.49\textwidth}
      \begin{pordentro}[Pular fases economiza tempo]
        Uma cor que anima só precisa de desenho; uma posição que muda, de layout e desenho. Ler o estado na fase mais tardia possível evita repetir as anteriores.
      \end{pordentro}
      \begin{dica}
        O \emph{Performance} do Flutter DevTools e o \emph{Layout Inspector} do Android Studio mostram onde o tempo do quadro foi gasto.
      \end{dica}
    \end{column}
  \end{columns}
\end{frame}

\begin{frame}{Por dentro: descrição, instância e renderização}
  \begin{columns}[T]
    \begin{column}{0.56\textwidth}
      \scriptsize
      \begin{tabularx}{\linewidth}{@{}L l l@{}}
        \toprule
        \cabecalho{Papel} & \cabecalho{Flutter} & \cabecalho{Compose} \\
        \midrule
        Descrição leve, recriada a cada mudança & \cod{Widget} & chamada \cod{@Composable} \\
        Instância duradoura, guarda o estado & \cod{Element} + \cod{State} & composição (\cod{remember}) \\
        Mede, posiciona e desenha & \cod{RenderObject} & \cod{LayoutNode} \\
        Informa leitores de tela & árvore de semântica & árvore de semântica \\
        \bottomrule
      \end{tabularx}
      \normalsize
      \begin{analogia}
        O widget é a planta do arquiteto: barata de redesenhar. O elemento é o apartamento construído, onde mora o estado. O objeto de renderização é a estrutura medida e pintada. Mudar a planta não obriga a demolir o prédio.
      \end{analogia}
    \end{column}
    \begin{column}{0.42\textwidth}
      \begin{conceito}[Por que separar]
        Descrições são criadas aos milhares por segundo, então precisam ser baratas. Medir e desenhar é caro, então os objetos de renderização são \textbf{reaproveitados} enquanto possível.
      \end{conceito}
      \begin{pordentro}[Semântica e acessibilidade]
        Em paralelo, o framework monta uma árvore de semântica (rótulos, papéis, ações) que o TalkBack e o VoiceOver leem. Ícone sem rótulo fica mudo para quem não enxerga (Módulo 04).
      \end{pordentro}
    \end{column}
  \end{columns}
\end{frame}

\begin{frame}{Atualizar sem refazer tudo}
  \begin{columns}[c]
    \begin{column}{0.42\textwidth}
      \centering
      \begin{tikzpicture}[level distance=10mm,
          every node/.style={caixa,font=\tiny,minimum height=5mm,minimum width=9mm,inner sep=2pt},
          edge from parent/.style={draw=azulpetroleo,line width=0.6pt},
          level 1/.style={sibling distance=22mm},
          level 2/.style={sibling distance=11mm}]
        \node {Tela}
          child {node {Barra} child {node {Título}}}
          child {node[caixaambarforte] (m) {Medidor}
            child {node[caixaambar] {Valor}}
            child {node[caixaambar] {Barra}}};
        \node[rotulo,below=11mm of m,text width=3.4cm] {o estado \cod{temperatura} mudou: só o \textbf{Medidor} e seus filhos são reconstruídos};
      \end{tikzpicture}
    \end{column}
    \begin{column}{0.56\textwidth}
      \begin{conceito}[Flutter: \emph{rebuild} e reconciliação]
        \cod{setState} marca o elemento como sujo; no próximo quadro, só ele executa \cod{build}. Filhos de mesmo tipo e chave são reaproveitados e só atualizados.
      \end{conceito}
      \begin{conceito}[Compose: recomposição inteligente]
        O Compose sabe quais funções leram cada estado e só executa essas de novo; as que recebem parâmetros iguais aos anteriores são puladas (\emph{skipping}).
      \end{conceito}
      \begin{atencao}[Consequência importante]
        Funções de interface podem rodar muitas vezes, em qualquer ordem: devem ser \textbf{rápidas e sem efeitos colaterais}. Rede e banco ficam em APIs de efeito (Módulo 07).
      \end{atencao}
    \end{column}
  \end{columns}
\end{frame}

\begin{frame}[fragile]{Identidade e chaves}
  \begin{columns}[T]
    \begin{column}{0.47\textwidth}
      \centering
      \begin{tikzpicture}[baseline=(current bounding box.north),every node/.style={font=\tiny}]
        \tikzset{item/.style={caixa,minimum width=1.5cm,minimum height=5mm,inner sep=1pt}}
        \node[rotulo] at (0.75,0.55) {antes};
        \node[item,caixaambar] (a) at (0.75,0) {A \quad expandido};
        \node[item] at (0.75,-0.6) {B};
        \node[item] at (0.75,-1.2) {C};
        \node[rotulo,text width=2cm] at (2.95,0.55) {remove A, \textbf{sem} chave};
        \node[item,caixaambar,draw=vermelhoerro] at (2.95,0) {B \quad expandido};
        \node[item] at (2.95,-0.6) {C};
        \node[rotulo,text width=2cm] at (5.15,0.55) {remove A, \textbf{com} chave};
        \node[item] at (5.15,0) {B};
        \node[item] at (5.15,-0.6) {C};
        \node[rotulo,text width=2cm,text=vermelhoerro] at (2.95,-1.25) {o estado ficou na posição 1};
        \node[rotulo,text width=2cm,text=verdeok] at (5.15,-1.25) {o estado saiu junto com A};
      \end{tikzpicture}
\begin{codigo}[title=Compose]{Kotlin}
items(itens, key = { it.id }) { eq ->
  ItemEquipamento(eq)
}
\end{codigo}
\begin{codigo}[title=Flutter]{Dart}
ItemEquipamento(
    key: ValueKey(eq.id), eq: eq, onTap: abrir)
\end{codigo}
    \end{column}
    \begin{column}{0.51\textwidth}
      \begin{definicao}[Chave (\emph{key})]
        Identificador estável que diz ao framework qual item é qual entre duas versões da árvore. Sem chave, a correspondência é feita pela \textbf{posição}.
      \end{definicao}
      \begin{conceito}[Quando importa]
        Em listas cujos itens têm estado próprio (expandido, campo digitado, animação) e podem ser inseridos, removidos ou reordenados.
      \end{conceito}
      \begin{dica}
        Use um identificador do domínio, como o \cod{id} do equipamento. Nunca o índice da lista nem um valor aleatório criado a cada reconstrução.
      \end{dica}
    \end{column}
  \end{columns}
\end{frame}

\begin{frame}{Temas, \emph{tokens} e contexto implícito}
  \begin{columns}[c]
    \begin{column}{0.42\textwidth}
      \centering
      \begin{tikzpicture}[node distance=7mm and 3mm,every node/.style={font=\tiny}]
        \node[caixaambarforte,minimum width=2.4cm,font=\tiny\bfseries] (t) {Tema};
        \node[caixa,below=of t,minimum width=1.6cm,font=\tiny] (s) {Tela};
        \node[caixa,below left=of s,minimum width=1.3cm,font=\tiny] (l) {Lista};
        \node[caixa,below right=of s,minimum width=1.3cm,font=\tiny] (b) {Botão};
        \node[caixa,below=of l,minimum width=1.3cm,font=\tiny] (c) {Cartão};
        \draw[seta] (t) -- (s); \draw[seta] (s) -- (l); \draw[seta] (s) -- (b); \draw[seta] (l) -- (c);
        \draw[setaambar,dashed] (t.east) -| (b.north);
        \draw[setaambar,dashed] (t.west) -| ([xshift=-3mm]c.west) -- (c.west);
        \node[rotulo,text width=4cm,below=2mm of c.south east] {tracejado: componentes leem o tema direto do contexto, sem receber parâmetros};
      \end{tikzpicture}
    \end{column}
    \begin{column}{0.56\textwidth}
      \begin{definicao}[\emph{Tokens} de design]
        Decisões visuais nomeadas pelo \textbf{papel}, não pelo valor: \cod{primary}, \cod{surface}, \cod{titleMedium}. O componente pede ``a cor primária''; o tema decide qual é.
      \end{definicao}
      \begin{pordentro}[Contexto implícito]
        Tema, idioma e tamanho da tela ficam disponíveis para toda uma subárvore: \cod{InheritedWidget} no Flutter (\cod{Theme.of(context)}), \cod{CompositionLocal} no Compose. Quem lê é atualizado quando o valor muda.
      \end{pordentro}
      \begin{dica}
        Com \emph{tokens}, tema escuro e cor dinâmica (Android 12+) funcionam sem mexer nos componentes.
      \end{dica}
    \end{column}
  \end{columns}
\end{frame}

\begin{frame}{Pixels, densidade e telas de todos os tamanhos}
  \begin{columns}[T]
    \begin{column}{0.49\textwidth}
      \centering
      \begin{tikzpicture}[baseline=(current bounding box.north)]
        \draw[cinzamedio,step=0.2] (0,0) grid (1.2,1.2);
        \fill[ambar,opacity=0.7] (0.2,0.2) rectangle (1.0,1.0);
        \node[rotulo] at (0.6,-0.25) {160 dpi: 48\,dp = 48\,px};
        \draw[cinzamedio,step=0.0667] (2.6,0) grid (3.8,1.2);
        \fill[ambar,opacity=0.7] (2.8,0.2) rectangle (3.6,1.0);
        \node[rotulo] at (3.2,-0.25) {480 dpi: 48\,dp = 144\,px};
      \end{tikzpicture}
      \raggedright
      \begin{definicao}[Pixel independente de densidade]
        \cod{px = dp} $\times$ \cod{(dpi} $\div$ \cod{160)}: um botão de 48\,dp tem o mesmo tamanho físico em qualquer tela. No Flutter, é o pixel lógico (fator em \cod{devicePixelRatio}).
      \end{definicao}
      \begin{dica}[Texto]
        Fontes também seguem a escala escolhida pelo usuário: \cod{sp} no Android, \cod{TextScaler} no Flutter.
      \end{dica}
    \end{column}
    \begin{column}{0.49\textwidth}
      \begin{conceito}[Responsivo $\times$ adaptativo]
        \textbf{Responsivo}: o mesmo layout se estica ou reorganiza. \textbf{Adaptativo}: um layout por classe de tamanho de janela (barra inferior no celular, trilho lateral no tablet). Pense em \textbf{espaço disponível}, não em modelo de aparelho.
      \end{conceito}
      \begin{atencao}[Android 16 e telas grandes]
        Com \cod{targetSdk} 36 (exigido pela Play desde agosto de 2026), telas com largura mínima $\geq$ 600\,dp ignoram travas de orientação e de redimensionamento; no API 37 não haverá exceção.
      \end{atencao}
    \end{column}
  \end{columns}
\end{frame}
